% lemma_stability.tex — stability of one hyperbola (second solver, 19.08): the orbit decomposition and the sharp constant
\subsection*{Stability: near-maximum sets are near a maximum set}

The classification of Theorem~\ref{thm:window} is in fact orbit-wise, and this makes the extremal problem for one hyperbola
\emph{stable}: a lawful set of size $3(p-1)-t$ differs from a maximum one in at most $t$ points.

\begin{lemma}[orbit decomposition]\label{lem:orbdec}
Let $O=\{\kappa,\sigma\kappa,\tau\kappa,\nu\kappa\}$ be a $V$-orbit of classes and $P_O\subseteq\Hc\cap\Gp$ the union of their lifts
($16$ points for a generic orbit, $8$ for an exceptional one).  Every line containing at least three points of $\Hc\cap\Gp$ has all of
them in a single $P_O$; consequently $\alpha(\Hc\cap\Gp)=\sum_O\alpha(P_O)$, and a set is lawful iff its restriction to every orbit is.
Moreover the incidence structure of $P_O$ --- which points lie on a common rich line --- is the same for all generic orbits and all $p$,
with rich-line profile $(4,4,3,3,3,3)$, and likewise for all exceptional orbits, with profile $(3,3)$.
\end{lemma}
\begin{proof}
The first assertion is part of Proposition~\ref{prop:lines}.  For the last one, a generic orbit consists of two classes of type $A$ and two
of type $D$, or of two of type $B$ and two of type $C$ (Lemma~\ref{lem:partners}: $\sigma$ preserves $A,D$ and exchanges $B,C$, $\tau$ the
opposite), and Proposition~\ref{prop:lines} then lists its rich lines: in the first case one four-point diagonal for the two $A$'s and one
for the two $D$'s (item (a)) and one three-point antidiagonal for each of the four classes (item (d)); in the second case the same with
diagonals and antidiagonals, and (b),(c) in place of (a),(d).  In both cases the incidence structure is the stated one, and the two cases
are exchanged by $(x,y)\mapsto(y,x)$.  An exceptional orbit is $\{\kappa,\nu\kappa\}$ with $\kappa$ fixed by $\sigma$ or by $\tau$, and
Proposition~\ref{prop:lines} gives two three-point lines.
\end{proof}

\begin{lemma}[a single orbit]\label{lem:orbopt}
For a generic orbit $\alpha(P_O)=12$ and the maximum set is \emph{unique}; for an exceptional orbit $\alpha(P_O)=6$ and there are exactly
nine maximum sets.  Moreover every lawful $S\subseteq P_O$ with $|S|=\alpha(P_O)-d$ satisfies $\min_M|S\setminus M|\le d$, the minimum
being over the maximum sets of $P_O$ (for an exceptional orbit $\min_M|S\setminus M|=0$: every lawful subset of size $6-d$ is contained in
a maximum set).
\end{lemma}
\begin{proof}
By Lemma~\ref{lem:orbdec} there are only two incidence structures, so this is a finite verification, carried out by exhaustive enumeration
of all subsets of the $16$, resp.\ $8$ points (\texttt{slack/\allowbreak t221/\allowbreak orbit\_\allowbreak stability.py}, \texttt{slack/\allowbreak t221/\allowbreak orbit\_\allowbreak types.py}; the values of
$\alpha(P_O)$, the number of maximum sets and the bound $\min_M|S\setminus M|\le d$ for $d\le3$ were confirmed for every orbit of every
prime $p\le31$, and exactly two incidence types occur; for $d\ge4$ the bound is trivial, since $|P_O\setminus M|=4$, resp.\ $2$).
The value $12$ and the count of maximum sets also follow from the gadget analysis of Theorem~\ref{thm:window}.
\end{proof}

\begin{proposition}[stability]\label{prop:stability}
Let $S\subseteq\Hc\cap\Gp$ be lawful with $|S|=3(p-1)-t$.  Then there is a maximum lawful set $M$,
$|M|=3(p-1)$, with $|S\setminus M|\le t$; in particular $|S\cap M|\ge3(p-1)-2t$.
\end{proposition}
\begin{proof}
By Lemma~\ref{lem:orbdec} the deficits $d_O:=\alpha(P_O)-|S\cap P_O|$ are non-negative and sum to $t$.  Choose in each orbit a maximum
set $M_O$ realising the minimum of Lemma~\ref{lem:orbopt}; the union $M=\bigcup_OM_O$ is lawful of size $\sum_O\alpha(P_O)=3(p-1)$
(Lemma~\ref{lem:orbdec}) and $|S\setminus M|=\sum_O|S\cap P_O\setminus M_O|\le\sum_Od_O=t$.
\end{proof}

\begin{remark}
(a) The bound is sharp for every $t$ in the near-maximum range: an exact computation (CP-SAT,
$p=11,13,17,19,23,29,31$ and $1\le t\le4$, every value proved optimal; \texttt{slack/t221/stability\_h3.py}) gives
\[
 \max\{\min_M|S\setminus M|\ :\ S\ \text{lawful},\ |S|=3(p-1)-t\}=\min(t,\,p-1-2s),
\]
the plateau value $p-1-2s$ being the number of generic classes: one unit of deficit buys exactly one deviating point until every
generic class has been touched.
(b) The proposition is the ``stability'' input that a proof of $\alpha(\Pset_k)\le3(p-1)+O(1)$ would want: it says that the part of a
lawful set on one hyperbola is, up to its own deficit, a piece of a HJSW-type configuration, so that the interaction with the second
hyperbola can be analysed against a fixed structure rather than against an arbitrary set.  What it does \emph{not} give is the exchange
rate: a point of the second hyperbola must be paid for by a deficit \emph{relative to the optimum} of the first, and that is the whole
conjecture.
(c) The same decomposition explains the count $9^s$ of Theorem~\ref{thm:window}: generic orbits are rigid, exceptional ones carry nine
patterns each.
\end{remark}
